\documentclass[reqno,11pt]{amsart}
\usepackage{math327-lecture}
\usepackage{tikz-cd}
\usepackage{needspace}
\lectureinfo{16}{Oct. 16, 2026}{Direct Products of Groups}

% Transcribed from handwritten Lecture 12, Sep. 29, 2025, pp. 1--10.
% The requested Sep. 28 source is filed and headed Sep. 29.
% Historical course announcements are omitted.
\begin{document}
\makelecturetitle

\section{Review: correspondence and the square}

Recall the surjective homomorphism $\varphi:S_4\to S_3$ from Lecture~15.
It sends a permutation of $\{1,2,3,4\}$ to its induced permutation of
the three partitions
\[
\begin{aligned}
  \Pi_1&=\bigl\{\{1,2\},\{3,4\}\bigr\},\\
  \Pi_2&=\bigl\{\{1,3\},\{2,4\}\bigr\},\\
  \Pi_3&=\bigl\{\{1,4\},\{2,3\}\bigr\}.
\end{aligned}
\]
Its kernel is
\[
  K=\{1,(1\,2)(3\,4),(1\,3)(2\,4),(1\,4)(2\,3)\}.
\]
The correspondence theorem identifies subgroups of $S_4$ containing
$K$ with subgroups of $S_3$, preserves normality, and gives
\[
  |\varphi^{-1}(\mathcal H)|=|K|\,|\mathcal H|=4|\mathcal H|.
\]
In particular, the square symmetries generated by the rotation
$(1\,2\,3\,4)$ and the diagonal reflection $(2\,4)$ form a subgroup
of order $8$ in $S_4$.

\section{External direct products}

There are two ways to approach products of groups:
\begin{itemize}
  \item An \emph{external} product starts with groups $G_1,G_2$ and
  constructs a new group $G_1\times G_2$.
  \item An \emph{internal} product starts with a group $G$ and asks
  whether suitable subgroups $G_1,G_2\leq G$ give an isomorphism
  $G\cong G_1\times G_2$.
\end{itemize}
Today we develop the external construction.

\begin{define}[External direct product]
Let $G_1$ and $G_2$ be groups. Their \emph{direct product}
$G_1\times G_2$ is the Cartesian product of their underlying sets,
with multiplication defined componentwise:
\[
  (g_1,g_2)(g'_1,g'_2)=(g_1g'_1,g_2g'_2).
\]
\end{define}

\begin{lemma}\label{lem:product-group}
The direct product is a group. Its identity is $(e_1,e_2)$ and
\[
  (g_1,g_2)^{-1}=(g_1^{-1},g_2^{-1}).
\]
\end{lemma}

\begin{exercise}
Verify the group axioms for $G_1\times G_2$, using the corresponding
axioms in each factor. If both groups are finite, show that
$|G_1\times G_2|=|G_1|\,|G_2|$.
\end{exercise}

\subsection{The factors as normal subgroups}

The maps
\[
  i_1:G_1\longrightarrow G_1\times G_2,\quad g_1\longmapsto(g_1,e_2),
\]
\[
  i_2:G_2\longrightarrow G_1\times G_2,\quad g_2\longmapsto(e_1,g_2)
\]
are injective homomorphisms. We often identify the factors with their
images, but the subgroups inside the product are, more precisely,
\[
  \widetilde G_1=G_1\times\{e_2\},\qquad
  \widetilde G_2=\{e_1\}\times G_2.
\]

\begin{lemma}
Both $\widetilde G_1$ and $\widetilde G_2$ are normal in
$G_1\times G_2$, and
\[
  \widetilde G_1\cap\widetilde G_2=\{(e_1,e_2)\}.
\]
\end{lemma}
\begin{proof}
For $(g_1,g_2)\in G_1\times G_2$ and $h_1\in G_1$,
\[
  (g_1,g_2)(h_1,e_2)(g_1,g_2)^{-1}
  =(g_1h_1g_1^{-1},e_2)\in\widetilde G_1.
\]
Thus $\widetilde G_1$ is normal; the same computation in the second
coordinate proves normality of $\widetilde G_2$.
An element in both subgroups must have first coordinate $e_1$ and
second coordinate $e_2$, proving the intersection statement.
\end{proof}

\section{The universal property of the product}

The coordinate projections
\[
  \pi_1:G_1\times G_2\to G_1,\quad (g_1,g_2)\mapsto g_1,
  \qquad
  \pi_2:G_1\times G_2\to G_2,\quad (g_1,g_2)\mapsto g_2
\]
are homomorphisms.

\begin{theorem}[Universal property]\label{thm:universal-product}
Given a group $G$ and homomorphisms
$\varphi:G\to G_1$ and $\psi:G\to G_2$, there is a unique
homomorphism $\Phi:G\to G_1\times G_2$ such that
\[
  \pi_1\circ\Phi=\varphi,\qquad \pi_2\circ\Phi=\psi.
\]
It is given by $\Phi(g)=(\varphi(g),\psi(g))$.
\end{theorem}

The following diagram records these identities:
\[
\begin{tikzcd}[column sep=large,row sep=large]
  G \arrow[r,dashed,"\Phi"] \arrow[rr,bend left=35,"\varphi"]
    \arrow[dr,"\psi"']
  & G_1\times G_2 \arrow[r,"\pi_1"] \arrow[d,"\pi_2"] & G_1\\
  & G_2 &
\end{tikzcd}
\]

\begin{proof}
Define $\Phi(g)=(\varphi(g),\psi(g))$. For $g,h\in G$,
\begin{align*}
  \Phi(gh)
  &=(\varphi(gh),\psi(gh))\\
  &=(\varphi(g)\varphi(h),\psi(g)\psi(h))\\
  &=(\varphi(g),\psi(g))(\varphi(h),\psi(h))\\
  &=\Phi(g)\Phi(h).
\end{align*}
Thus $\Phi$ is a homomorphism, and its projections are $\varphi$ and
$\psi$. Conversely, those projections specify both coordinates of
$\Phi(g)$, so they determine $\Phi$ uniquely.
\end{proof}

Equivalently, there is a bijection of sets
\[
  \Hom(G,G_1\times G_2)
  \longleftrightarrow\Hom(G,G_1)\times\Hom(G,G_2),
\]
sending $\Phi$ to $(\pi_1\circ\Phi,\pi_2\circ\Phi)$, with inverse
$(\varphi,\psi)\mapsto[g\mapsto(\varphi(g),\psi(g))]$.

\section{Products of cyclic groups}

Write $C_d$ for a cyclic group of order $d$.

\begin{lemma}\label{lem:coprime-cyclic}
If $r,s\geq1$ and $\gcd(r,s)=1$, then
\[
  C_{rs}\cong C_r\times C_s.
\]
\end{lemma}

\begin{warningbox}
The coprimality assumption matters. For example,
$C_4\not\cong C_2\times C_2$. Every element $(u,v)$ of
$C_2\times C_2$ satisfies
\[
  (u,v)^2=(u^2,v^2)=(1,1),
\]
whereas $C_4$ has an element of order $4$. An isomorphism preserves
element orders, so these groups cannot be isomorphic.
\end{warningbox}

\begin{proof}[Proof of Lemma~\ref{lem:coprime-cyclic}]
Choose generators
\[
  C_{rs}=\langle x\rangle,\qquad C_r=\langle z\rangle,
  \qquad C_s=\langle y\rangle.
\]
Define
\[
  \varphi:C_{rs}\to C_r,\quad x^i\mapsto z^i,
  \qquad
  \psi:C_{rs}\to C_s,\quad x^i\mapsto y^i.
\]
These maps are well defined: if $x^i=x^j$, then $rs\mid(i-j)$,
so both $r$ and $s$ divide $i-j$, giving $z^i=z^j$ and $y^i=y^j$.
They are homomorphisms, since, for instance,
\[
  \varphi(x^ix^j)=z^{i+j}=z^iz^j.
\]
By the universal property, they combine into the homomorphism
\[
  \Phi:C_{rs}\longrightarrow C_r\times C_s,
  \qquad x^i\longmapsto(z^i,y^i).
\]
The two coordinate kernels are
\[
  \ker(\varphi)=\langle x^r\rangle,\qquad
  \ker(\psi)=\langle x^s\rangle.
\]
Consequently, $x^i\in\ker(\Phi)$ if and only if $r\mid i$ and
$s\mid i$. Since $\gcd(r,s)=1$, this is equivalent to $rs\mid i$,
so $\ker(\Phi)=\{1\}$ and $\Phi$ is injective.

Both its domain and its codomain have order $rs$. Thus $\Phi$ is also
surjective, and hence an isomorphism. Equivalently, the counting
formula gives $|\im(\Phi)|=rs$, so its image is all of
$C_r\times C_s$.
\end{proof}

\end{document}
